\documentclass[11pt,paper=a4,answers]{exam}
\usepackage{graphicx,lastpage}
\usepackage{upgreek}
\usepackage{censor}
\usepackage{gensymb}
\usepackage{amsmath}
\usepackage{multicol}
\usepackage{enumitem}
\usepackage{tasks}
\usepackage{adjustbox}
\usepackage{xcolor}
\usepackage{tfrupee}
\setlength{\voffset}{0.25in}
\usepackage{tabularx}
\newcolumntype{Y}{>{\centering\arraybackslash}X}
%==============================================================
% Customise Non-Essential Packages Below
% =============================================================
\usepackage[most]{tcolorbox}
\usepackage{eso-pic}
\usepackage{tikz}
\AddToShipoutPictureBG{%
\begin{tikzpicture}[remember picture,overlay]
\foreach \x in {-8,-4,0,4,8,12,16,20}
\foreach \y in {-8,-4,0,4,8,12,16,20} {
\node[opacity=0.05,rotate=45,anchor=center] at ([xshift=\x cm,yshift=\y cm]current page.center) {%
\begin{minipage}{2cm}
\centering
\includegraphics[width=2cm]{images/logo_black.png}\\
\small preppeo.com
\end{minipage}%
};
}
\end{tikzpicture}%
}
\printanswers

%==============================================================
% Customise Non-Essential Packages Above
% =============================================================
\definecolor{svgray}{rgb}{0.90, 0.90, 0.90}
\graphicspath{{images/}}

\usepackage[normalem]{ulem}
\renewcommand\ULthickness{1pt}   %%---> For changing thickness of underline
\setlength\ULdepth{3.5ex}%\maxdimen ---> For changing depth of underline
\renewcommand{\baselinestretch}{1}
\chead{
\includegraphics[scale=0.1,height=2cm ]{logo.png}
}
\pagestyle{empty}

\pagestyle{headandfoot}
\headrule
\cfoot{W: https://preppeo.com; M: +91-8130711689}

\begin{document}
%==============================================================
% Customise Question paper details below this
% =============================================================
\begin{center}
    \centering
    \renewcommand{\arraystretch}{1.5}
    \begin{tabularx}{\textwidth}{YYY}
    & \normalsize\bf{{\Large{{Quantitative Reasoning}}}} & \\
    By Preppeo & \normalsize\bf{{sat\_lid\_024}} & SAT\\
    \normalsize{{\textbf{{Unit:}} Advanced Math}} & \normalsize{{\textbf{{Chapter:}} Exponential functions and equations}} & \normalsize{{\textbf{{Topic:}} Compound Interest \& Population Growth}} \\
    \end{tabularx}
\end{center}
\par\noindent\rule{\textwidth}{1pt}
% =============================================================
% Customise Question paper details above this
% =============================================================

\begin{questions}
\pointsinrightmargin
\pointsdroppedatright
\marksnotpoints
\pointpoints{mark}{marks}
\pointformat{\boldmath\themarginpoints}
\bracketedpoints

% =============================================================
% Question Begin
% =============================================================
% Lid Topic: Advanced Math — Compound Interest & Population Growth
% Total Questions: 50

% --- PART 1: Core Growth Modeling (1-25) ---

% Question 1
% Category: Concept | Difficulty: Medium | Question Type: Multiple Choice
\question
An investment account was opened with an initial deposit of $1,000$. The account earns 5\% interest compounded annually. Which of the following functions models the value $V(t)$, in dollars, of the account $t$ years after the initial deposit?
\begin{tasks}(2)
    \task $V(t) = 1,000(0.05)^t$
    \task $V(t) = 1,000(1.05)^t$
    \task $V(t) = 1,000 + 0.05t$
    \task $V(t) = 1,000(1.5)^t$
\end{tasks}

\begin{solution}
\textbf{Conceptual Explanation:} \\
Compound interest follows the exponential growth model $y = a(1 + r)^t$. The "base" of the exponent must be the growth factor, which is $1$ plus the interest rate expressed as a decimal.

\vspace{20pt}

\textbf{Calculation and Logic:} \\
Initial deposit ($a$) = 1,000. \\
Interest rate ($r$) = 5\% = 0.05. \\
Growth factor = $1 + 0.05 = 1.05$.

\vspace{10pt}

Substituting these into the model: \\
$V(t) = 1,000(1.05)^t$.

\vspace{25pt}

Answer: \colorbox{yellow!100}{B}
\end{solution}

\vspace{40pt}

% Question 2
% Category: Exam level | Difficulty: Hard | Question Type: Numeric Entry
\question
The population of a certain town is estimated to increase by 6\% every year. If the population was 25,000 in 2020, what is the estimated population of the town in 2022?

\begin{solution}
\textbf{Conceptual Explanation:} \\
For multi-year growth, apply the growth factor repeatedly for each year that passes.

\vspace{20pt}

\textbf{Calculation and Logic:} \\
Initial population = 25,000. \\
Time elapsed = $2022 - 2020 = 2$ years. \\
Growth factor = $1 + 0.06 = 1.06$.

\vspace{10pt}

Population in 2022 = $25,000 \times (1.06)^2$. \\
$25,000 \times 1.1236 = 28,090$.

\vspace{25pt}

Answer: \colorbox{yellow!100}{28090}
\end{solution}

\vspace{40pt}

% Question 3
% Category: Practice | Difficulty: Medium | Question Type: Multiple Choice
\question
$P(t) = 4,500(1.035)^t$ \\
The function $P$ models the population of a colony of insects $t$ weeks after an initial observation. Which of the following is the best interpretation of the value 1.035 in the function?
\begin{tasks}(1)
    \task The population increases by 3.5 each week.
    \task The population increases by 3.5\% each week.
    \task The population increases by 103.5\% each week.
    \task The initial population was 1,035.
\end{tasks}

\begin{solution}
\textbf{Conceptual Explanation:} \\
In the form $y = a(b)^t$, the base $b$ represents the growth factor. If $b > 1$, the percentage growth rate is $(b - 1) \times 100$.

\vspace{20pt}

\textbf{Calculation and Logic:} \\
Base = 1.035. \\
Rate = $1.035 - 1 = 0.035$. \\
Convert to percentage: $0.035 \times 100 = 3.5\%$.

\vspace{25pt}

Answer: \colorbox{yellow!100}{B}
\end{solution}

\vspace{40pt}

% Question 4
% Category: Concept | Difficulty: Medium | Question Type: Multiple Choice
\question
An amount of $P$ dollars is deposited into an account that pays $r\%$ annual interest compounded annually. If no additional deposits or withdrawals are made, which of the following represents the total amount in the account after 10 years?
\begin{tasks}(2)
    \task $P(1 + r)^{10}$
    \task $P(1 + \frac{r}{100})^{10}$
    \task $P + 10r$
    \task $P(r)^{10}$
\end{tasks}

\begin{solution}
\textbf{Conceptual Explanation:} \\
When an interest rate is given as a percentage "r", it must be divided by 100 to be used as a decimal in an exponential growth formula.

\vspace{25pt}

Answer: \colorbox{yellow!100}{B}
\end{solution}

\vspace{40pt}

% Question 5
% Category: Exam level | Difficulty: Hard | Question Type: Numeric Entry
\question
$A = 500(1.02)^{4t}$ \\
The equation above models the amount of money $A$, in dollars, in a savings account $t$ years after the initial deposit. The account earns interest compounded $n$ times per year. What is the value of $n$?

\begin{solution}
\textbf{Conceptual Explanation:} \\
The standard compound interest formula is $A = P(1 + \frac{r}{n})^{nt}$. The coefficient of $t$ in the exponent represents the number of compounding periods per year ($n$).

\vspace{20pt}

\textbf{Calculation and Logic:} \\
In the given model $A = 500(1.02)^{4t}$, the exponent is $4t$. \\
Comparing this to $nt$, we see that $n = 4$.

\vspace{25pt}

Answer: \colorbox{yellow!100}{4}
\end{solution}

\vspace{40pt}

% Question 6
% Category: Practice | Difficulty: Medium | Question Type: Multiple Choice
\question
A population of 100,000 decreases by 2\% each year. Which function $P(t)$ models the population $t$ years from now?
\begin{tasks}(2)
    \task $P(t) = 100,000(1.02)^t$
    \task $P(t) = 100,000(0.98)^t$
    \task $P(t) = 100,000(0.02)^t$
    \task $P(t) = 100,000 - 0.02t$
\end{tasks}

\begin{solution}
\textbf{Conceptual Explanation:} \\
Exponential decay uses a base of $(1 - r)$. A 2\% decrease results in a multiplier of $1 - 0.02 = 0.98$.

\vspace{25pt}

Answer: \colorbox{yellow!100}{B}
\end{solution}

\vspace{40pt}

% Question 7
% Category: Exam level | Difficulty: Hard | Question Type: Numeric Entry
\question
For the population function $f$, the table below shows the population $x$ years after a survey was conducted. If $f(x) = a(b)^x$, what is the value of $b$?
\begin{center}
\begin{tabular}{|c|c|}
\hline
$x$ & $f(x)$ \\ \hline
0 & 8,000 \\ \hline
1 & 8,240 \\ \hline
2 & 8,487.2 \\ \hline
\end{tabular}
\end{center}

\begin{solution}
\textbf{Conceptual Explanation:} \\
In an exponential function, the base $b$ is the ratio between any two consecutive output values: $b = \frac{f(x+1)}{f(x)}$.

\vspace{20pt}

\textbf{Calculation and Logic:} \\
Use the values for $x=0$ and $x=1$: \\
$b = \frac{8,240}{8,000}$ \\
$b = \frac{824}{800} = 1.03$

\vspace{25pt}

Answer: \colorbox{yellow!100}{1.03}
\end{solution}

\vspace{40pt}

% Question 8
% Category: Concept | Difficulty: Medium | Question Type: Multiple Choice
\question
The function $M(t) = 15,000(1.045)^t$ models the amount of money in a retirement account $t$ years after 2015. What was the amount of money in the account in 2015?
\begin{tasks}(4)
    \task $1.045$
    \task $15,000$
    \task $15,675$
    \task $675$
\end{tasks}

\begin{solution}
\textbf{Conceptual Explanation:} \\
The value in 2015 represents the "initial value" ($t = 0$). In the form $y = a(b)^t$, the initial value is the constant $a$.

\vspace{25pt}

Answer: \colorbox{yellow!100}{B}
\end{solution}

\vspace{40pt}

% Question 9
% Category: Exam level | Difficulty: Hard | Question Type: Multiple Choice
\question
The population of Lowell was 90,000 in 2010 and has been increasing by 4\% each year. Which of the following functions best models the population of Lowell $n$ decades after 2010? (1 decade = 10 years)
\begin{tasks}(2)
    \task $p(n) = 90,000(1.04)^{10n}$
    \task $p(n) = 90,000(1.04)^{\frac{n}{10}}$
    \task $p(n) = 90,000(1.40)^n$
    \task $p(n) = 90,000(1.04n)^{10}$
\end{tasks}

\begin{solution}
\textbf{Conceptual Explanation:} \\
If $t$ is years and $n$ is decades, then $t = 10n$. Substitute $10n$ for the time variable in the standard growth formula.

\vspace{20pt}

\textbf{Calculation and Logic:} \\
Standard yearly model: $P = 90,000(1.04)^t$. \\
Since $t = 10n$, the decade model is: \\
$P = 90,000(1.04)^{10n}$.

\vspace{25pt}

Answer: \colorbox{yellow!100}{A}
\end{solution}

\vspace{40pt}

% Question 10
% Category: Practice | Difficulty: Hard | Question Type: Numeric Entry
\question
An investment of \$2,000 grows at an annual rate of 10\% compounded annually. What will be the value of the investment, in dollars, after 3 years?

\begin{solution}
\textbf{Calculation and Logic:} \\
Year 1: $2,000 \times 1.1 = 2,200$. \\
Year 2: $2,200 \times 1.1 = 2,420$. \\
Year 3: $2,420 \times 1.1 = 2,662$.

\vspace{25pt}

Answer: \colorbox{yellow!100}{2662}
\end{solution}

% Question 11
% Category: Concept | Difficulty: Medium | Question Type: Multiple Choice
\question
The population $P$ of a bacteria culture doubles every 3 hours. If the initial population is $N$, which of the following functions models the population after $t$ hours?
\begin{tasks}(2)
    \task $P(t) = N(2)^{3t}$
    \task $P(t) = N(2)^{\frac{t}{3}}$
    \task $P(t) = N(3)^{2t}$
    \task $P(t) = N(3)^{\frac{t}{2}}$
\end{tasks}

\begin{solution}
\textbf{Conceptual Explanation:} \\
In a doubling model $y = a(2)^{\frac{t}{d}}$, the variable $d$ is the "doubling time." This ensures that when $t=d$, the exponent is 1 and the population has doubled once.

\vspace{25pt}

Answer: \colorbox{yellow!100}{B}
\end{solution}

\vspace{40pt}

% Question 12
% Category: Exam level | Difficulty: Hard | Question Type: Numeric Entry
\question
$f(t) = 300(1.21)^{\frac{t}{2}}$ \\
The function $f$ models the number of users on a website $t$ months after it launched. The function can be rewritten in the form $f(t) = 300(1 + \frac{r}{100})^t$. What is the value of $r$?

\begin{solution}
\textbf{Conceptual Explanation:} \\
Use the power of a power property $(a^m)^n = a^{mn}$ to combine the fractional exponent with the base. This converts the growth factor to a monthly basis.

\vspace{20pt}

\textbf{Calculation and Logic:} \\
$(1.21)^{\frac{t}{2}} = (1.21^{1/2})^t$. \\
Calculate the square root: $\sqrt{1.21} = 1.1$. \\
The monthly growth factor is 1.1. \\
$1 + \frac{r}{100} = 1.1 \Rightarrow \frac{r}{100} = 0.1 \Rightarrow r = 10$.

\vspace{25pt}

Answer: \colorbox{yellow!100}{10}
\end{solution}

\vspace{40pt}

% Question 13
% Category: Practice | Difficulty: Medium | Question Type: Multiple Choice
\question
A town's population increases by 100\% every 15 years. If the population is currently 5,000, what will it be in 45 years?
\begin{tasks}(4)
    \task 15,000
    \task 20,000
    \task 40,000
    \task 80,000
\end{tasks}

\begin{solution}
\textbf{Conceptual Explanation:} \\
A 100\% increase is the same as doubling. Determine how many 15-year periods occur in 45 years.

\vspace{20pt}

\textbf{Calculation and Logic:} \\
Cycles = $45 / 15 = 3$ doublings. \\
$5,000 \times 2 = 10,000$ (15 yrs) \\
$10,000 \times 2 = 20,000$ (30 yrs) \\
$20,000 \times 2 = 40,000$ (45 yrs).

\vspace{25pt}

Answer: \colorbox{yellow!100}{C}
\end{solution}

\vspace{40pt}

% Question 14
% Category: Exam level | Difficulty: Hard | Question Type: Numeric Entry
\question
The value of an investment increases by $x\%$ each year. If the value of the investment triples every 12 years, what is the value of $(1 + \frac{x}{100})^{36}$?

\begin{solution}
\textbf{Conceptual Explanation:} \\
The expression $(1 + \frac{x}{100})^{36}$ represents the total growth factor over 36 years. Determine how many 12-year tripling periods fit into 36 years.

\vspace{20pt}

\textbf{Calculation and Logic:} \\
Number of periods = $36 / 12 = 3$. \\
Since the value triples each period, the total growth multiplier is $3 \times 3 \times 3 = 3^3 = 27$.

\vspace{25pt}

Answer: \colorbox{yellow!100}{27}
\end{solution}

\vspace{40pt}

% Question 15
% Category: Concept | Difficulty: Medium | Question Type: Multiple Choice
\question
A savings account offers an annual interest rate of 4\% compounded semi-annually. Which expression represents the value of a \$500 deposit after $t$ years?
\begin{tasks}(2)
    \task $500(1.04)^t$
    \task $500(1.02)^{2t}$
    \task $500(1.04)^{2t}$
    \task $500(1.08)^t$
\end{tasks}

\begin{solution}
\textbf{Conceptual Explanation:} \\
"Semi-annually" means $n = 2$. Divide the annual rate by 2 and multiply the time $t$ by 2.

\vspace{20pt}

\textbf{Calculation and Logic:} \\
Periodic rate = $4\% / 2 = 2\% = 0.02$. \\
Compounding periods = $2t$. \\
$V = 500(1 + 0.02)^{2t} = 500(1.02)^{2t}$.

\vspace{25pt}

Answer: \colorbox{yellow!100}{B}
\end{solution}

\vspace{40pt}

% Question 16
% Category: Exam level | Difficulty: Hard | Question Type: Numeric Entry
\question
A population of rare plants is decreasing by 10\% each year. If there are 1,000 plants now, after how many years will the population be 729?

\begin{solution}
\textbf{Calculation and Logic:} \\
Year 1: $1,000 \times 0.9 = 900$. \\
Year 2: $900 \times 0.9 = 810$. \\
Year 3: $810 \times 0.9 = 729$.

\vspace{25pt}

Answer: \colorbox{yellow!100}{3}
\end{solution}

\vspace{40pt}

% Question 17
% Category: Practice | Difficulty: Medium | Question Type: Multiple Choice
\question
Which of the following describes a population that grows by 50\% every year?
\begin{tasks}(1)
    \task The population increases by 50 individuals each year.
    \task The population in any year is 1.5 times the population of the previous year.
    \task The population in any year is 0.5 times the population of the previous year.
    \task The population triples every 2 years.
\end{tasks}

\begin{solution}
\textbf{Conceptual Explanation:} \\
A 50\% increase corresponds to a growth factor of $1 + 0.50 = 1.5$. This means the new total is 1.5 times the previous total.

\vspace{25pt}

Answer: \colorbox{yellow!100}{B}
\end{solution}

\vspace{40pt}

% Question 18
% Category: Exam level | Difficulty: Hard | Question Type: Numeric Entry
\question
$V(t) = 1,200(1.02)^t$ \\
The function above models the value of an account $t$ years after 2000. By what percentage does the value of the account increase every 5 years? (Round to the nearest percent).

\begin{solution}
\textbf{Calculation and Logic:} \\
Growth factor for 1 year = 1.02. \\
Growth factor for 5 years = $(1.02)^5 \approx 1.104$. \\
Rate = $1.104 - 1 = 0.104 \approx 10\%$.

\vspace{25pt}

Answer: \colorbox{yellow!100}{10}
\end{solution}

\vspace{40pt}

% Question 19
% Category: Concept | Difficulty: Medium | Question Type: Multiple Choice
\question

Which of the following equations could represent the graph above?
\begin{tasks}(2)
    \task $y = 50(0.8)^x$
    \task $y = 50(1.2)^x$
    \task $y = 50 - 1.2x$
    \task $y = 50 + 1.2x$
\end{tasks}

\begin{solution}
\textbf{Conceptual Explanation:} \\
A curve that starts above the x-axis and bends upward represents exponential growth, where the base must be greater than 1.

\vspace{25pt}

Answer: \colorbox{yellow!100}{B}
\end{solution}

\vspace{40pt}

% Question 20
% Category: Exam level | Difficulty: Hard | Question Type: Numeric Entry
\question
The population of a city doubles every 20 years. If the population was 40,000 in 1980, what was the population in 1940?

\begin{solution}
\textbf{Conceptual Explanation:} \\
Work backward in time. Each 20-year period into the past represents a halving of the population.

\vspace{20pt}

\textbf{Calculation and Logic:} \\
Time difference = $1980 - 1940 = 40$ years. \\
Number of 20-year periods = $40 / 20 = 2$. \\
Population in 1960 = $40,000 / 2 = 20,000$. \\
Population in 1940 = $20,000 / 2 = 10,000$.

\vspace{25pt}

Answer: \colorbox{yellow!100}{10000}
\end{solution}

\vspace{40pt}

% Question 21
% Category: Practice | Difficulty: Medium | Question Type: Multiple Choice
\question
A bank account pays 1\% annual interest compounded monthly. Which of the following is the monthly interest rate?
\begin{tasks}(4)
    \task 0.01
    \task 0.12
    \task 0.000833
    \task 0.00833
\end{tasks}

\begin{solution}
\textbf{Calculation and Logic:} \\
Annual rate = 1\% = 0.01. \\
Monthly rate = $0.01 / 12 \approx 0.000833$.

\vspace{25pt}

Answer: \colorbox{yellow!100}{C}
\end{solution}

\vspace{40pt}

% Question 22
% Category: Concept | Difficulty: Medium | Question Type: Multiple Choice
\question
Which table shows exponential growth?
\begin{tasks}(2)
    \task \begin{tabular}{|c|c|c|c|} \hline $x$ & 1 & 2 & 3 \\ \hline $y$ & 10 & 20 & 30 \\ \hline \end{tabular}
    \task \begin{tabular}{|c|c|c|c|} \hline $x$ & 1 & 2 & 3 \\ \hline $y$ & 10 & 20 & 40 \\ \hline \end{tabular}
    \task \begin{tabular}{|c|c|c|c|} \hline $x$ & 1 & 2 & 3 \\ \hline $y$ & 10 & 5 & 2.5 \\ \hline \end{tabular}
    \task \begin{tabular}{|c|c|c|c|} \hline $x$ & 1 & 2 & 3 \\ \hline $y$ & 10 & 12 & 14 \\ \hline \end{tabular}
\end{tasks}

\begin{solution}
\textbf{Conceptual Explanation:} \\
Growth requires output values to increase by a common ratio. In Choice B, the ratio is $20/10 = 2$ and $40/20 = 2$.

\vspace{25pt}

Answer: \colorbox{yellow!100}{B}
\end{solution}

\vspace{40pt}

% Question 23
% Category: Exam level | Difficulty: Hard | Question Type: Numeric Entry
\question
$125^x = 5^{12}$ \\
What is the value of $x$?

\begin{solution}
\textbf{Calculation and Logic:} \\
$(5^3)^x = 5^{12}$. \\
$3x = 12 \Rightarrow x = 4$.

\vspace{25pt}

Answer: \colorbox{yellow!100}{4}
\end{solution}

\vspace{40pt}

% Question 24
% Category: Practice | Difficulty: Hard | Question Type: Multiple Choice
\question
The population of a town triples every 8 years. If the current population is $P$, what will the population be after 24 years?
\begin{tasks}(4)
    \task $3P$
    \task $9P$
    \task $24P$
    \task $27P$
\end{tasks}

\begin{solution}
\textbf{Calculation and Logic:} \\
Periods = $24 / 8 = 3$. \\
Multiplier = $3^3 = 27$.

\vspace{25pt}

Answer: \colorbox{yellow!100}{D}
\end{solution}

\vspace{40pt}

% Question 25
% Category: Exam level | Difficulty: Hard | Question Type: Numeric Entry
\question
A radioactive isotope has a half-life of 2 days. If a sample has 64 grams initially, how many grams will remain after 10 days?

\begin{solution}
\textbf{Calculation and Logic:} \\
Cycles = $10 / 2 = 5$. \\
$64 \times (0.5)^5 = 64 \times 0.03125 = 2$.

\vspace{25pt}

Answer: \colorbox{yellow!100}{2}
\end{solution}

% --- PART 2: Advanced Compound Interest & Population (26-50) ---

% Question 26
% Category: Exam level | Difficulty: Hard | Question Type: Numeric Entry
\question
An investment account earns 10\% annual interest compounded semi-annually. If the initial deposit is \$1,000, what is the value of the account, in dollars, after 1 year?

\begin{solution}
\textbf{Conceptual Explanation:} \\
Semi-annual compounding means the interest is applied twice a year ($n=2$). You must divide the annual rate by 2 and apply it for 2 periods.

\vspace{20pt}

\textbf{Calculation and Logic:} \\
Periodic rate = $10\% / 2 = 5\% = 0.05$. \\
Amount after 1st period (6 months) = $1,000 \times 1.05 = 1,050$. \\
Amount after 2nd period (12 months) = $1,050 \times 1.05 = 1,102.50$.

\vspace{25pt}

Answer: \colorbox{yellow!100}{1102.5}
\end{solution}

\vspace{40pt}

% Question 27
% Category: Practice | Difficulty: Hard | Question Type: Multiple Choice
\question
A certain population grows according to the model $P(t) = P_0 (1.04)^t$, where $t$ is the number of years. Which of the following best models the population $d$ decades after the start? (1 decade = 10 years)
\begin{tasks}(2)
    \task $P(d) = P_0 (1.04)^{\frac{d}{10}}$
    \task $P(d) = P_0 (1.04)^{10d}$
    \task $P(d) = P_0 (1.40)^d$
    \task $P(d) = P_0 (1.04)^{d}$
\end{tasks}

\begin{solution}
\textbf{Conceptual Explanation:} \\
Substitute the relationship between the time units into the exponent. Since $t = 10d$, the exponent $t$ becomes $10d$.

\vspace{25pt}

Answer: \colorbox{yellow!100}{B}
\end{solution}

\vspace{40pt}

% Question 28
% Category: Exam level | Difficulty: Hard | Question Type: Numeric Entry
\question
The population of a city was 50,000 in 2010. If the population grew at a constant annual rate of 2\% and reached 60,950 after $k$ years, what is the value of $k$?

\begin{solution}
\textbf{Calculation and Logic:} \\
$50,000(1.02)^k = 60,950$ \\
$(1.02)^k = 1.219$ \\
Testing integer values: \\
$1.02^5 \approx 1.104$ \\
$1.02^{10} \approx 1.219$.

\vspace{25pt}

Answer: \colorbox{yellow!100}{10}
\end{solution}

\vspace{40pt}

% Question 29
% Category: Practice | Difficulty: Medium | Question Type: Numeric Entry
\question
What is the growth factor for an investment that increases by 12.5\% each year?

\begin{solution}
\textbf{Calculation and Logic:} \\
Growth factor = $1 + r = 1 + 0.125 = 1.125$.

\vspace{25pt}

Answer: \colorbox{yellow!100}{1.125}
\end{solution}

\vspace{40pt}

% Question 30
% Category: Concept | Difficulty: Medium | Question Type: Multiple Choice
\question
If the population of a town is modeled by $P(t) = 20,000(0.85)^t$, what is the annual percentage decrease?
\begin{tasks}(4)
    \task 0.15\%
    \task 1.5\%
    \task 15\%
    \task 85\%
\end{tasks}

\begin{solution}
\textbf{Calculation and Logic:} \\
$0.85 = 1 - r \Rightarrow r = 0.15 = 15\%$.

\vspace{25pt}

Answer: \colorbox{yellow!100}{C}
\end{solution}

\vspace{40pt}

% Question 31
% Category: Exam level | Difficulty: Hard | Question Type: Numeric Entry
\question
An account with an initial balance of \$5,000 earns 8\% annual interest compounded quarterly. What is the interest rate applied at the end of each quarter as a decimal?

\begin{solution}
\textbf{Calculation and Logic:} \\
Quarterly means $n = 4$. \\
Periodic rate = $0.08 / 4 = 0.02$.

\vspace{25pt}

Answer: \colorbox{yellow!100}{0.02}
\end{solution}

\vspace{40pt}

% Question 32
% Category: Practice | Difficulty: Hard | Question Type: Numeric Entry
\question
The value of a certain car depreciates by 15\% each year. If the car is worth \$20,000 today, how much will it be worth, in dollars, in 2 years?

\begin{solution}
\textbf{Calculation and Logic:} \\
$20,000 \times (0.85)^2 = 20,000 \times 0.7225 = 14,450$.

\vspace{25pt}

Answer: \colorbox{yellow!100}{14450}
\end{solution}

\vspace{40pt}

% Question 33
% Category: Exam level | Difficulty: Hard | Question Type: Multiple Choice
\question
$f(t) = 100(1.05)^t$ \\
$g(t) = 100(1.10)^{\frac{t}{2}}$ \\
Which of the following statements comparing $f$ and $g$ is true for $t > 0$?
\begin{tasks}(1)
    \task Function $g$ grows at a faster rate than function $f$.
    \task Function $f$ grows at a faster rate than function $g$.
    \task Both functions grow at the same rate.
    \task The initial value of $g$ is greater than the initial value of $f$.
\end{tasks}

\begin{solution}
\textbf{Conceptual Explanation:} \\
Convert $g(t)$ to have a standard $t$ exponent: \\
$(1.10)^{1/2} \approx 1.0488$. \\
Since $1.05 > 1.0488$, function $f$ has a slightly higher annual growth rate.

\vspace{25pt}

Answer: \colorbox{yellow!100}{B}
\end{solution}

\vspace{40pt}

% Question 34
% Category: Practice | Difficulty: Medium | Question Type: Numeric Entry
\question
A population triples every 5 years. By what factor does the population increase in 15 years?

\begin{solution}
\textbf{Calculation and Logic:} \\
Periods = $15 / 5 = 3$. \\
Factor = $3^3 = 27$.

\vspace{25pt}

Answer: \colorbox{yellow!100}{27}
\end{solution}

\vspace{40pt}

% Question 35
% Category: Exam level | Difficulty: Hard | Question Type: Numeric Entry
\question
If $5^{x+2} = 125$, what is the value of $x$?

\begin{solution}
\textbf{Calculation and Logic:} \\
$5^{x+2} = 5^3$. \\
$x + 2 = 3 \Rightarrow x = 1$.

\vspace{25pt}

Answer: \colorbox{yellow!100}{1}
\end{solution}

\vspace{40pt}

% Question 36
% Category: Practice | Difficulty: Hard | Question Type: Multiple Choice
\question
An investment of $P$ dollars increases in value by 20\% every 4 years. Which function models the value $V(t)$ after $t$ years?
\begin{tasks}(2)
    \task $V(t) = P(1.20)^{4t}$
    \task $V(t) = P(1.20)^{\frac{t}{4}}$
    \task $V(t) = P(0.20)^{\frac{t}{4}}$
    \task $V(t) = P(1.05)^t$
\end{tasks}

\begin{solution}
\textbf{Conceptual Explanation:} \\
The rate 20\% gives a base of 1.20. The exponent must represent the number of 4-year cycles, which is $t/4$.

\vspace{25pt}

Answer: \colorbox{yellow!100}{B}
\end{solution}

\vspace{40pt}

% Question 37
% Category: Practice | Difficulty: Medium | Question Type: Numeric Entry
\question
If a bank account with \$1,000 earns 4\% interest compounded annually, how much interest, in dollars, is earned in the first year?

\begin{solution}
\textbf{Calculation and Logic:} \\
$1,000 \times 0.04 = 40$.

\vspace{25pt}

Answer: \colorbox{yellow!100}{40}
\end{solution}

\vspace{40pt}

% Question 38
% Category: Exam level | Difficulty: Hard | Question Type: Numeric Entry
\question
A town's population was 10,000 in 2000 and 20,000 in 2010. If the population continues to double every 10 years, what will the population be in 2040?

\begin{solution}
\textbf{Calculation and Logic:} \\
2010 to 2040 = 30 years (3 doublings). \\
$20,000 \times 2 = 40,000$ (2020) \\
$40,000 \times 2 = 80,000$ (2030) \\
$80,000 \times 2 = 160,000$ (2040).

\vspace{25pt}

Answer: \colorbox{yellow!100}{160000}
\end{solution}

\vspace{40pt}

% Question 39
% Category: Concept | Difficulty: Medium | Question Type: Multiple Choice
\question
Which of the following represents an exponential function where the initial value is 5 and the growth rate is 100\%?
\begin{tasks}(4)
    \task $y = 5(1)^x$
    \task $y = 5(2)^x$
    \task $y = 100(5)^x$
    \task $y = 5 + 100x$
\end{tasks}

\begin{solution}
\textbf{Calculation and Logic:} \\
Initial ($a$) = 5. \\
Growth factor = $1 + 1.00 = 2$.

\vspace{25pt}

Answer: \colorbox{yellow!100}{B}
\end{solution}

\vspace{40pt}

% Question 40
% Category: Exam level | Difficulty: Hard | Question Type: Numeric Entry
\question
$2 \cdot 4^x = 32$ \\
What is the value of $x$?

\begin{solution}
\textbf{Calculation and Logic:} \\
$4^x = 16$. \\
$4^2 = 16 \Rightarrow x = 2$.

\vspace{25pt}

Answer: \colorbox{yellow!100}{2}
\end{solution}

\vspace{40pt}

% Question 41
% Category: Practice | Difficulty: Hard | Question Type: Numeric Entry
\question
A savings account starts with \$1,000. It earns 2\% interest compounded annually. How much more money, in dollars, will be in the account after 2 years compared to after 1 year?

\begin{solution}
\textbf{Calculation and Logic:} \\
Year 1: $1,000 \times 1.02 = 1,020$. \\
Year 2: $1,020 \times 1.02 = 1,040.40$. \\
Difference = $1,040.40 - 1,020 = 20.40$.

\vspace{25pt}

Answer: \colorbox{yellow!100}{20.4}
\end{solution}

\vspace{40pt}

% Question 42
% Category: Exam level | Difficulty: Hard | Question Type: Multiple Choice
\question
A substance has a half-life of 8 hours. What fraction of the original amount remains after 24 hours?
\begin{tasks}(4)
    \task 1/3
    \task 1/4
    \task 1/8
    \task 1/16
\end{tasks}

\begin{solution}
\textbf{Calculation and Logic:} \\
Cycles = $24 / 8 = 3$. \\
$(1/2)^3 = 1/8$.

\vspace{25pt}

Answer: \colorbox{yellow!100}{C}
\end{solution}

\vspace{40pt}

% Question 43
% Category: Practice | Difficulty: Hard | Question Type: Numeric Entry
\question
$f(n) = 10 \cdot 2^n$ \\
$g(n) = 10 \cdot 3^n$ \\
What is the value of $\frac{g(2)}{f(2)}$?

\begin{solution}
\textbf{Calculation and Logic:} \\
$g(2) = 10 \times 9 = 90$. \\
$f(2) = 10 \times 4 = 40$. \\
Ratio = $90 / 40 = 2.25$.

\vspace{25pt}

Answer: \colorbox{yellow!100}{2.25}
\end{solution}

\vspace{40pt}

% Question 44
% Category: Concept | Difficulty: Medium | Question Type: Multiple Choice
\question
In the compound interest formula $A = P(1 + \frac{r}{n})^{nt}$, what does $P$ represent?
\begin{tasks}(1)
    \task The annual interest rate
    \task The number of times interest is compounded per year
    \task The principal (initial amount)
    \task The total time in years
\end{tasks}

\begin{solution}
\textbf{Conceptual Explanation:} \\
$P$ stands for Principal, the starting amount of money.

\vspace{25pt}

Answer: \colorbox{yellow!100}{C}
\end{solution}

\vspace{40pt}

% Question 45
% Category: Exam level | Difficulty: Hard | Question Type: Numeric Entry
\question
If a population increases by 100\% every decade, what is the annual growth factor rounded to two decimal places?

\begin{solution}
\textbf{Calculation and Logic:} \\
$r^{10} = 2$. \\
$r = 2^{1/10} \approx 1.0717$. \\
Growth factor $\approx 1.07$.

\vspace{25pt}

Answer: \colorbox{yellow!100}{1.07}
\end{solution}

\vspace{40pt}

% Question 46
% Category: Practice | Difficulty: Medium | Question Type: Numeric Entry
\question
A population starts at 500 and grows by 10\% each year. What is the population after 1 year?

\begin{solution}
\textbf{Calculation and Logic:} \\
$500 \times 1.10 = 550$.

\vspace{25pt}

Answer: \colorbox{yellow!100}{550}
\end{solution}

\vspace{40pt}

% Question 47
% Category: Exam level | Difficulty: Hard | Question Type: Multiple Choice
\question
An amount is deposited into an account that triples every 10 years. After how many years will the amount be 9 times the initial deposit?
\begin{tasks}(4)
    \task 18
    \task 20
    \task 30
    \task 90
\end{tasks}

\begin{solution}
\textbf{Calculation and Logic:} \\
Since $9 = 3^2$, it takes 2 tripling periods. \\
$2 \times 10 = 20$ years.

\vspace{25pt}

Answer: \colorbox{yellow!100}{B}
\end{solution}

\vspace{40pt}

% Question 48
% Category: Practice | Difficulty: Hard | Question Type: Numeric Entry
\question
If $(1.03)^x = 1.092727$, what is the value of $x$ to the nearest integer?

\begin{solution}
\textbf{Calculation and Logic:} \\
$1.03^2 = 1.0609$. \\
$1.03^3 = 1.092727$.

\vspace{25pt}

Answer: \colorbox{yellow!100}{3}
\end{solution}

\vspace{40pt}

% Question 49
% Category: Exam level | Difficulty: Hard | Question Type: Numeric Entry
\question
A town's population decreases by 50\% every 20 years. If the population is 8,000 now, what was it 40 years ago?

\begin{solution}
\textbf{Calculation and Logic:} \\
Backward 20 years: $8,000 \times 2 = 16,000$. \\
Backward 40 years: $16,000 \times 2 = 32,000$.

\vspace{25pt}

Answer: \colorbox{yellow!100}{32000}
\end{solution}

\vspace{40pt}

% Question 50
% Category: Practice | Difficulty: Medium | Question Type: Numeric Entry
\question
An initial investment of \$1,000 earns 5\% annual interest compounded annually. What is the value of the account after 2 years?

\begin{solution}
\textbf{Calculation and Logic:} \\
$1,000 \times (1.05)^2 = 1,000 \times 1.1025 = 1,102.50$.

\vspace{25pt}

Answer: \colorbox{yellow!100}{1102.5}
\end{solution}

% =============================================================
% Question End
% =============================================================

\end{questions}
\begin{center}
\rule{.5\textwidth}{1pt}
\end{center}
\end{document}
